\subsection{Angular Coordinate Interface (ACI)}
\label{subsec:aci}
Motivated by the classifier geometry mismatch identified in our
analysis, we introduce an Angular Coordinate Interface (ACI) that
separates the student's directional class relations from
magnitude- and bias-dependent score modulation.
ACI modifies only the student classifier; the pre-trained teacher
remains frozen and retains its original classification head.

\subsubsection{Cosine Classifier}
A conventional affine classifier combines angular similarity
with feature magnitudes, class-weight norms, and additive biases.
To remove these factors from the student readout, we replace
its affine head with a bias-free cosine classifier.
Given a student feature $\mathbf{h}$ and the weight vector
$\mathbf{w}_k$ for class $k$, the angular logit is
\begin{equation}
    \tilde{z}_k
    = s \cdot
    \frac{\mathbf{w}_k^\top \mathbf{h}}
    {\|\mathbf{w}_k\|_2 \|\mathbf{h}\|_2}
    = s \cos\theta_k,
    \label{eq:aci_logit}
\end{equation}
where $\theta_k$ is the angle between the feature and the
class-weight vector, and $s>0$ is a fixed scale shared across
samples and classes.

Feature normalization removes sample-dependent magnitude
scaling, while class-weight normalization removes
class-specific multiplicative factors. Omitting the bias
further eliminates additive class-specific offsets.
Unlike a global temperature adjustment, ACI therefore changes
the score parameterization rather than merely rescaling
the resulting logits.

For nonzero vectors, the resulting scores are invariant to
positive rescaling of the feature or individual class weights,
while both feature and class-weight directions remain learnable.
Thus, ACI makes the student scores depend only on angular
relations and the shared scale, without requiring direct
teacher--student feature alignment.
These scores support both label-supervised learning and the
subsequent SICD objective: ACI defines the student output
interface, whereas SICD specifies the distillation signal.
